\section{Related Work}

\paragraph{Long-term Dialogue.} 
Recent approaches involve retrieving historical context from a range of previous dialogues and reasoning over retrieved segments in a temporal order~\cite{lee2023prompted, lu2023memochat, zhong2023memorybank, liang2023unleashing} and/or using events to scaffold the dialogues \cite{jang2023conversation, zhang2023mind} to enable consistency in long-term conversations. Some limitations of such frameworks are:
(1) The accuracy of retrieval can be compromised, as the retrieval model is generally trained on tasks focusing on semantic similarity rather than specifically on such dialogues.
Additionally, real-world dialogues often feature co-references and missing content (\textit{i.e.,} anaphora)~\cite{anantha2021open}, which further complicate the retrieval process~\cite{mallen-etal-2023-trust, gao2023enabling, liu2023evaluating};
(2) Challenges arise in reasoning over retrieved documents, especially when the model struggles to identify the correct context among the retrieved data~\cite{liu2023lost};
(3) Reasoning over time intervals presents challenges. For example, the way a system responds about past events can vary depending on the amount of time that has passed since the last conversation~\cite{zhang2023mind, jang2023conversation}.
Therefore, it is essential to have conversations of considerable length, as well as a systematic evaluation framework, to accurately assess the effectiveness of approaches to long-term dialogue generation. We design a long-term conversation generation pipeline based on retrieval augmentation and events graphs and propose a framework for evaluating long-term dialog agents.

\paragraph{Multi-modal Dialogue.}
Multi-modal dialogue primarily consists of two types of tasks: image-grounded dialogue and image-sharing dialogue. 
The image-grounded dialogue task is centered around responding to questions~\cite{antol2015vqa, das2017visual, kottur2019clevr} or creating natural conversations related to specific images~\cite{mostafazadeh2017image, shuster2020image, meng2020openvidial, zheng2021mmchat}. 
Conversely, the image-sharing dialogue task focuses on selecting images that semantically align with the provided dialogue context~\cite{zang2021photochat, feng2023mmdialog, lee2023dialogcc}. We use a method from the image-sharing dialogue task to create multimodal dialogs which are then evaluated as an image-grounded dialogue task. 


\paragraph{Synthetic Evaluation Benchmark.}
Faced with a shortage of human-generated data and observing that LLMs are approaching the quality of human-level annotations~\cite{he2023annollm, lee-etal-2023-making}, there has been a surge in research drawing inspiration from this development. 
Consequently, numerous studies have started utilizing LLMs to augment or synthesize large-scale dialogue benchmarks for assessing responses in everyday social interactions~\cite{kim-etal-2023-soda}, examining responses in multi-modal environment~\cite{feng2023mmdialog}, and evaluating responses that align with specific persona~\cite{jandaghi2023faithful}. We leverage LLMs to create data but ensure its high quality with human verification and editing.
